\section{Methodological Considerations and Model Comparisons}

\subsection{Landmark Analysis and Longitudinal Data}
Longitudinal data present unique challenges and opportunities for predictive modeling, particularly in clinical settings. 
Landmark analysis has emerged as a robust methodology for managing time-dependent covariates and outcomes, facilitating dynamic 
predictions at various clinically relevant points (landmarks). In our study, we employed landmark analysis explicitly, 
defining each clinical visit as a landmark to predict mortality within 90 days following each landmark.

\subsection{Rationale for Dynamic Predictions}
Initially, we trained separate predictive models at each landmark using an XGBoost classifier, each time utilizing clinical 
data available strictly up to that landmark. This explicit separation ensured methodological rigor, prevented data leakage, 
and provided clear insights into how predictive performance evolved as additional patient history was accumulated.

\subsection{Single vs. Multiple Models for Landmark Analysis}
While initially training separate XGBoost models at each landmark, we considered whether a single, unified model trained once
on the entire landmark dataset could yield better predictive performance. This approach aligns closely with real-world 
clinical scenarios, where clinicians typically rely on a single continuously used predictive model rather than repeatedly trained separate models.

To explore this hypothesis rigorously, we adapted our modeling strategy to train a single XGBoost model across all landmarks 
and then dynamically applied this model at subsequent landmarks. Explicit precautions were taken to ensure that each landmark prediction 
only included data available up to that point, thereby preventing data leakage.

\subsection{Application and Adaptation of Large Language Models (LLMs)}
Parallel to traditional machine learning methods, we investigated the applicability of transformer-based large language models (LLMs), 
specifically MedBERT, to leverage narrative patient clinical histories. Initially, we fine-tuned separate MedBERT models at each landmark, 
observing moderate performance (AUC $\sim$0.8). We hypothesized that the limited performance might result from repeatedly fine-tuning separate 
models on smaller subsets of data, thus limiting the model's ability to leverage full patient trajectories.

To address this, we thought and implemented a single-model fine-tuning strategy. We fine-tuned one MedBERT model on the combined landmark dataset, 
which allowed the model to exploit richer and more comprehensive narrative clinical contexts. The subsequent dynamic evaluation at each landmark 
provided improved consistency and potentially higher predictive accuracy due to increased training data and the continuity of clinical narrative context. 
However this introduce data leakage, since in this way the model implicitly sees future information because:
\begin{itemize}
    \item When the model is learning from Landmark 1 data, it also sees Landmark 2 and Landmark 3 data for the same patient during training.
    \item Landmark 2 and Landmark 3 narratives explicitly include future information (visits) relative to Landmark 1.
\end{itemize}

\subsection{Narrative Prompt Design for MedBERT}
We recognized that the quality and clarity of the input provided to MedBERT significantly impact predictive performance. 
Initially, our prompts consisted of simple concatenations of diagnoses, medications, and procedures. However, realizing that 
transformer-based models excel at interpreting coherent narratives, we refined our prompts into structured, clinically realistic 
narratives explicitly describing patient demographics, medical histories, medications, and procedural histories clearly and cohesively. 
This change substantially aligns the input format more closely with MedBERT's pre-training corpus, thereby potentially enhancing predictive performance.

\subsection{Scientific Rigor and Data Leakage}
A critical methodological point addressed in our study was the explicit prevention of data leakage. We emphasized that each landmark's 
predictive inputs must strictly contain information available at or before the landmark, never including future information. 
By explicitly structuring our datasets to meet this criterion, we maintained high methodological rigor and ensured scientific validity.

\subsection{Comparative Performance Insights}
Interestingly, despite methodological improvements, our experiments demonstrated that the XGBoost model leveraging structured numeric features 
(such as age, intervals between visits, and structured clinical variables) consistently outperformed MedBERT. The substantial improvement in 
predictive performance observed at higher landmarks for XGBoost highlights the strength of structured, numeric features in longitudinal clinical predictions.

\subsection{Conclusions and Recommendations}
Our findings indicate that while LLMs such as MedBERT offer promising avenues for leveraging rich narrative clinical data, traditional 
structured machine learning methods like XGBoost remain particularly effective for longitudinal predictive modeling tasks involving structured 
numeric data. Future research should further explore hybrid approaches or enhance narrative input structures for LLMs to better exploit their 
inherent strengths. Explicit comparisons and methodological transparency remain crucial to guide model selection in clinical predictive modeling tasks.

